\documentclass[11pt]{scrartcl}

\usepackage[top=1.5cm]{geometry}
\usepackage{url}
\usepackage{float}
\usepackage{listings}
\usepackage{xcolor}
\usepackage{graphicx}

\setlength{\parindent}{0em}
\setlength{\parskip}{0.5em}

\newcommand{\youranswerhere}{[Your answer goes here \ldots]}
\renewcommand{\thesubsection}{\arabic{subsection}}

\lstdefinestyle{dmrsql}{
  language=SQL,
  basicstyle=\small\ttfamily,
  keywordstyle=\color{magenta!75!black},
  stringstyle=\color{green!50!black},
  showspaces=false,
  showstringspaces=false,
  commentstyle=\color{gray}}

\lstdefinestyle{dmrJava}{
  language=JAVA,
  basicstyle=\small\ttfamily,
  keywordstyle=\color{magenta!75!black},
  stringstyle=\color{green!50!black},
  showspaces=false,
  showstringspaces=false,
  commentstyle=\color{gray}}

\title{
  \textbf{\large Projektaufgabe 2 } \\
  Phase 2 – Legen der Basis (2.5 P) \\
  {\large Datenmanagement jenseits von Relationen}
}

\author{
  Gruppen Nummer (A17) \\
  \large Acikyol Aleyna, 12032843 \\
  \large Özcan Ahmet, 12311287 \\
  \large Sadykova Nargiz, 12129337 \\
  
}

\begin{document}

\maketitle\thispagestyle{empty}

Dieses Reporting Template dient der Vorbereitung der Abgabe von Phase 2.

\subsection*{Datengenerator für Matrizen mit Sparsity (0.5 Punkte)}

Zeigen Sie den Code der Funktion generate() als Listing oder Screenshot. Gehen Sie
(kurz) auf die wesentlichen Aspekte ein

{Die Funktion generate(l, sparsity) erzeugt zwei nahezu quadratische Matrizen A und B als 2D-Arrays mit Gleitkommazahlen. Aus dem Größenparameter l werden die Dimensionen m = l - 1 und n = l - 1 berechnet, sodass A die Größe (l-1) × l und B die Größe l × (l-1) besitzt. Die Matrixeinträge werden zunächst mit zufälligen double-Werten erzeugt; in Python entsprechen die verwendeten float-Werte der double-Genauigkeit. Anschließend wird für jede Matrix aus der Gesamtanzahl der Zellen und dem Parameter sparsity die Anzahl der Nullwerte bestimmt. Diese Nullwerte werden zufällig auf die möglichen Matrixpositionen verteilt, indem alle Indexpunkte (i, j) erzeugt, gemischt und die berechnete Anzahl davon auf 0.0 gesetzt wird. Die Funktion gibt beide Matrizen als Rückgabewert zurück.}

\begin{lstlisting}[style=dmrsql]
[import random

def generate(l, sparsity):
    if l < 2:
        raise ValueError("l muss mindestens 2 sein")
    if not (0.0 <= sparsity <= 1.0):
        raise ValueError("sparsity muss zwischen 0 und 1 liegen")

    m = l - 1
    n = l - 1

    # A: m x l
    A = [[random.uniform(1.0, 10.0) for _ in range(l)] for _ in range(m)]
    # B: l x n
    B = [[random.uniform(1.0, 10.0) for _ in range(n)] for _ in range(l)]

    total_A = m * l
    total_B = l * n

    zeros_in_A = round(sparsity * total_A)
    zeros_in_B = round(sparsity * total_B)

    positions_A = [(i, j) for i in range(m) for j in range(l)]
    positions_B = [(i, j) for i in range(l) for j in range(n)]

    random.shuffle(positions_A)
    random.shuffle(positions_B)

    for i, j in positions_A[:zeros_in_A]:
        A[i][j] = 0.0
    for i, j in positions_B[:zeros_in_B]:
        B[i][j] = 0.0

    return A, B
]
\end{lstlisting}

\subsection*{Import der Matrizen in das DBMS (0.5 Punkte):}

Geben Sie das Create Table Statement für Matrix A an.

\begin{lstlisting}[style=dmrsql]
[CREATE TABLE A (
    i INT,
    j INT,
    val DOUBLE PRECISION
);
import psycopg2
from psycopg2.extras import execute_values

def import_to_db(matrix, table_name, db_config):
    # Daten vorbereiten
    sparse_data = []
    for i in range(len(matrix)):
        for j in range(len(matrix[0])):
            if matrix[i][j] != 0.0:
                sparse_data.append((i + 1, j + 1, matrix[i][j]))
    
    # Verbindung und Import
    conn = psycopg2.connect(**db_config)
    cursor = conn.cursor()
    
    query = f"INSERT INTO {table_name} (i, j, val) VALUES %s"
    execute_values(cursor, query, sparse_data)
    
    conn.commit()
    cursor.close()
    conn.close()
    print(f"Import in {table_name} erfolgreich!")]
\end{lstlisting}

Für die Übung bereiten Sie eine Demo des Datenimports vor.


\subsection*{Wahl des Toy Beispiels (0.5 P):}

Geben Sie Matrix A und B als 2D Array an.

\youranswerhere{}

Zeigen Sie die äquivalente darstellung der Matrix A und B in der Datenbank

\youranswerhere{}

\subsection*{Implementierung von Ansatz 0 (0.5 P):}
Zeigen Sie den Code der Matrixmultiplikation als Listing oder Screenshot. Erläutern Sie
(kurz), welche Laufzeit ihr Algorithmus hat und warum das Kriterium keinen Algorith-
mus mit sub-kubischer Laufzeit zu w¨ahlen erf¨ullt ist

\begin{lstlisting}[style=dmrsql]
[Matrixmultiplikation ...]
\end{lstlisting}

\subsection*{Implementierung von Ansatz 1 (0.5 P):}
Berechnen Sie von Hand das Ergebnis $C = A \times B$ füur ihr Toy Beispiel und geben Sie es
nachfolgend an.

\youranswerhere{}

Zeigen Sie, dass Ihr System $C$ korrekt berechnet (z.B. als Screenshot)

\subsection*{Zeitmamagement}

Benötigte Zeit pro Person (nur Phase 1): \textbf{Aleyna Acikyol: ,Ahmet Özcan: ,Nargiz Sadykova: 3.5 Stunden.}

\subsection*{References}

\begin{table}[H]
  \centering
  \begin{tabular}{c}
    \hline
    \textbf{Important:} Reference your information sources! \tabularnewline
    Remove this section if you use footnotes to reference your information sources. \tabularnewline
    \hline
  \end{tabular}
\end{table}

\end{document}
